% GENERATED FILE. Edit canonical lessons, then run npm run generate:books.
% canonical-source-hash: fnv1a64:edb158fd22323a05
% canonical-lessons: ML-C339-ayirunnu, ML-C339-vayikkukayayirunnu, ML-C339-pokumayirunnu, ML-R339-the-past-three-ways

\chapter{Was, Was Doing, Used To}
\label{ch:ml-was-was-doing-used-to}
\hlchaptermodality{339}

\begin{chapteropening}
\textbf{By the end of this chapter:} \emph{I can say where someone was and what there was, set a scene with what was going on, and say what used to happen.}

Complete ML-R339-the-past-three-ways: was and there was, a scene in the past progressive with an event cutting into it, and a habit with -umāyirunnu.
\end{chapteropening}

\section[\texorpdfstring{\ml{ആയിരുന്നു} (āyirunnu)}{ആയിരുന്നു (āyirunnu)}]{\ml{ആയിരുന്നു} (āyirunnu) — was, and there was}

\label{lesson:ML-C339-ayirunnu}



\begin{quote}
Say \textquotedblleft{}he is at home\textquotedblright{}: \textbf{\ml{അവൻ} \ml{വീട്ടിൽ} \ml{ആണ്}} (\emph{avan vīṭṭil āṇŭ}). Then \textquotedblleft{}there is a small shop here\textquotedblright{}: \textbf{\ml{ഇവിടെ} \ml{ഒരു} \ml{ചെറിയ} \ml{കട} \ml{ഉണ്ട്}} (\emph{iviṭe oru ceṟiya kaṭa uṇṭŭ}).

A story needs both of these in the past.
\end{quote}

\subsection*{You'll want to know: \ml{ആയിരുന്നു}}
\begin{quote}
\textbf{\ml{അവൻ} \ml{ഇന്നലെ} \ml{വീട്ടിൽ} \ml{ആയിരുന്നു}.}
\end{quote}

\begin{quote}
\emph{avan innale vīṭṭil āyirunnu.} — \textquotedblleft{}He was at home yesterday.\textquotedblright{}
\end{quote}

Where \textbf{\ml{ആണ്}} says \textquotedblleft{}is\textquotedblright{}, \textbf{\ml{ആയിരുന്നു}} says \textquotedblleft{}was\textquotedblright{}. It does not change for the person: \emph{ñān āyirunnu}, \emph{nī āyirunnu}, \emph{avar āyirunnu}: I was, you were, they were.

\begin{quote}
\textbf{\ml{അത്} \ml{ഒരു} \ml{നല്ല} \ml{സിനിമ} \ml{ആയിരുന്നു}.}
\end{quote}

\begin{quote}
\emph{atŭ oru nalla sinima āyirunnu.} — \textquotedblleft{}That was a good film.\textquotedblright{}
\end{quote}

In print it is often joined to the word before it: \textbf{\ml{വീട്ടിലായിരുന്നു}}.

\subsection*{You'll want to know: \ml{ഉണ്ടായിരുന്നു}}
Where \textbf{\ml{ഉണ്ട്}} says \textquotedblleft{}there is\textquotedblright{} or, with the dative, \textquotedblleft{}somebody has\textquotedblright{}, \textbf{\ml{ഉണ്ടായിരുന്നു}} says \textquotedblleft{}there was\textquotedblright{} and \textquotedblleft{}somebody had\textquotedblright{}.

\begin{quote}
\textbf{\ml{ഗ്രാമത്തിൽ} \ml{ഒരു} \ml{ചെറിയ} \ml{കട} \ml{ഉണ്ടായിരുന്നു}.}
\end{quote}

\begin{quote}
\emph{grāmattil oru ceṟiya kaṭa uṇṭāyirunnu.} — \textquotedblleft{}There was a small shop in the village.\textquotedblright{}
\end{quote}

\begin{quote}
\textbf{\ml{മുത്തശ്ശിക്ക്} \ml{ഒരു} \ml{വലിയ} \ml{വീട്} \ml{ഉണ്ടായിരുന്നു}.}
\end{quote}

\begin{quote}
\emph{muttaśśikkŭ oru valiya vīṭŭ uṇṭāyirunnu.} — \textquotedblleft{}Grandmother had a big house.\textquotedblright{}
\end{quote}

The negative follows the same road: \textbf{\ml{ഇല്ല}} becomes \textbf{\ml{ഇല്ലായിരുന്നു}} (\emph{illāyirunnu}), \textquotedblleft{}there was not\textquotedblright{}.

\begin{quote}
\textbf{\ml{അവിടെ} \ml{ഫോൺ} \ml{ഇല്ലായിരുന്നു}.}
\end{quote}

\begin{quote}
\emph{aviṭe phōṇ illāyirunnu.} — \textquotedblleft{}There was no phone there.\textquotedblright{}
\end{quote}

\begin{grammarlens}[title={Grammar Lens: two pieces you half know}]
\textbf{\ml{ആയിരുന്നു}} is two pieces: \textbf{\ml{ആയി}}, \textquotedblleft{}having become\textquotedblright{}, and \textbf{\ml{ഇരുന്നു}}, the past of the verb \textbf{\ml{ഇരിക്കുക}}, \textquotedblleft{}to sit, to stay\textquotedblright{}. So \emph{avan vīṭṭil āyirunnu} is, piece by piece, \textquotedblleft{}he, at home, having become, stayed\textquotedblright{}. \textbf{\ml{ഉണ്ടായിരുന്നു}} is \textbf{\ml{ഉണ്ട്}} with the same \textbf{\ml{ആയിരുന്നു}} behind it.

So the past keeps the present's split: \textbf{\ml{ആയിരുന്നു}} for what or where, \textbf{\ml{ഉണ്ടായിരുന്നു}} for being there at all, or having.
\end{grammarlens}

\subsection*{Your turn}
Say these aloud:

\begin{itemize}
\raggedright
  \item \textquotedblleft{}He was at home yesterday\textquotedblright{} — \emph{avan innale vīṭṭil āyirunnu}
  \item \textquotedblleft{}That was a good film\textquotedblright{} — \emph{atŭ oru nalla sinima āyirunnu}
  \item \textquotedblleft{}Grandmother had a big house\textquotedblright{} — \emph{muttaśśikkŭ oru valiya vīṭŭ uṇṭāyirunnu}
  \item \textquotedblleft{}There was no phone there\textquotedblright{} — \emph{aviṭe phōṇ illāyirunnu}
\end{itemize}

\begin{tcolorbox}[breakable,colback=teal!4,colframe=teal!35!black,title={Before you move on}]
\textquotedblleft{}Was\textquotedblright{}, as in \textquotedblleft{}he was at home\textquotedblright{}? (\textbf{\ml{ആയിരുന്നു}}, \emph{āyirunnu}.) \textquotedblleft{}There was\textquotedblright{}, or \textquotedblleft{}had\textquotedblright{}? (\textbf{\ml{ഉണ്ടായിരുന്നു}}, \emph{uṇṭāyirunnu}.) \textquotedblleft{}There was not\textquotedblright{}? (\textbf{\ml{ഇല്ലായിരുന്നു}}, \emph{illāyirunnu}.)
\end{tcolorbox}
\section[\texorpdfstring{\ml{വായിക്കുകയായിരുന്നു} (vāyikkukayāyirunnu)}{വായിക്കുകയായിരുന്നു (vāyikkukayāyirunnu)}]{\ml{വായിക്കുകയായിരുന്നു} (vāyikkukayāyirunnu) — was reading: the scene and the event}

\label{lesson:ML-C339-vayikkukayayirunnu}



\begin{quote}
Say \textquotedblleft{}I read\textquotedblright{} for today, then for yesterday. (\textbf{\ml{ഞാൻ} \ml{വായിക്കുന്നു}} (\emph{ñān vāyikkunnu}), \textbf{\ml{ഞാൻ} \ml{വായിച്ചു}} (\emph{ñān vāyiccu}).) Then the word for \textquotedblleft{}was\textquotedblright{}. (\emph{āyirunnu}.)

\emph{ñān vāyiccu} has an edge: it happened, and it is over. A story needs a past with no edge.
\end{quote}

\begin{grammarlens}[title={Grammar Lens: the dictionary form, then \ml{ആയിരുന്നു}}]
Take the dictionary form of the verb, the one ending in \textbf{-\ml{ഉക}}, and put \textbf{\ml{ആയിരുന്നു}} after it. A \textbf{\ml{യ}} slides in where the two vowels meet:

\begin{quote}
\ml{വായിക്കുക} + \ml{ആയിരുന്നു} $\to$ \textbf{\ml{വായിക്കുകയായിരുന്നു}} — \emph{vāyikkukayāyirunnu} — was reading
\end{quote}

So \textbf{\ml{നടക്കുകയായിരുന്നു}} (\emph{naṭakkukayāyirunnu}), \textquotedblleft{}was walking\textquotedblright{}, and \textbf{\ml{ഉറങ്ങുകയായിരുന്നു}} (\emph{uṟaṅṅukayāyirunnu}), \textquotedblleft{}was sleeping\textquotedblright{}. It is one word for every person. Quick speech often shortens it to \emph{vāyikkuvāyirunnu}; the full form is the written one.
\end{grammarlens}

\begin{grammarlens}[title={Grammar Lens: the scene, and what cut into it}]
A told story splits its verbs in two. The \textbf{-\ml{ഉകയായിരുന്നു}} past paints the \textbf{scene}: what was going on. The plain past drops an \textbf{event} into it: what happened.

\begin{quote}
\textbf{\ml{ഞാൻ} \ml{പുസ്തകം} \ml{വായിക്കുകയായിരുന്നു}. \ml{അപ്പോൾ} \ml{അമ്മ} \ml{വിളിച്ചു}.}
\end{quote}

\begin{quote}
\emph{ñān pustakaṁ vāyikkukayāyirunnu. appōḷ amma viḷiccu.} — \textquotedblleft{}I was reading a book. Then Mother called.\textquotedblright{}
\end{quote}

The reading was still running when the call cut into it. The call is over in a moment. \textbf{\ml{അപ്പോൾ}}, \textquotedblleft{}at that moment\textquotedblright{}, is the joint between them.

\begin{quote}
\textbf{\ml{ഞങ്ങൾ} \ml{കടപ്പുറത്ത്} \ml{നടക്കുകയായിരുന്നു}. \ml{പെട്ടെന്ന്} \ml{മഴ} \ml{പെയ്തു}.}
\end{quote}

\begin{quote}
\emph{ñaṅṅaḷ kaṭappuṟattŭ naṭakkukayāyirunnu. peṭṭennŭ maḻa peytu.} — \textquotedblleft{}We were walking on the beach. Suddenly it rained.\textquotedblright{}
\end{quote}

\textbf{\ml{മഴ}} is \textquotedblleft{}rain\textquotedblright{}, and \textbf{\ml{പെയ്തു}} is what rain does: it \textquotedblleft{}fell\textquotedblright{}. Again, the walk is the scene, and the rain is the event.
\end{grammarlens}

\subsection*{Your turn}
Say these aloud:

\begin{itemize}
\raggedright
  \item \textquotedblleft{}I was reading\textquotedblright{}, \textquotedblleft{}I was walking\textquotedblright{}, \textquotedblleft{}I was sleeping\textquotedblright{}
  \item \emph{ñān pustakaṁ vāyikkukayāyirunnu. appōḷ amma viḷiccu.}
  \item the walk on the beach, then the rain that cut into it
  \item which verb in each pair is the scene, and which the event
\end{itemize}

\begin{tcolorbox}[breakable,colback=teal!4,colframe=teal!35!black,title={Before you move on}]
How do you build \textquotedblleft{}was reading\textquotedblright{}? (\textbf{\ml{വായിക്കുക}} + \textbf{\ml{ആയിരുന്നു}}: \emph{\textbf{vāyikkukayāyirunnu}}.) In the reading sentence, which verb is the scene? (\emph{\textbf{vāyikkukayāyirunnu}}, still running.) Which is the event? (\emph{\textbf{viḷiccu}}, over in a moment.)
\end{tcolorbox}
\section[\texorpdfstring{\ml{പോകുമായിരുന്നു} (pōkumāyirunnu)}{പോകുമായിരുന്നു (pōkumāyirunnu)}]{\ml{പോകുമായിരുന്നു} (pōkumāyirunnu) — used to go}

\label{lesson:ML-C339-pokumayirunnu}



\begin{quote}
Say \textquotedblleft{}tomorrow I will go\textquotedblright{}. (\textbf{\ml{നാളെ} \ml{ഞാൻ} \ml{പോകും}} (\emph{nāḷe ñān pōkuṁ}).) Then \textquotedblleft{}I was going\textquotedblright{}. (\emph{ñān pōkukayāyirunnu}.)

The first has the \textbf{-\ml{ഉം}} of the future. Today it moves into the past.
\end{quote}

\begin{grammarlens}[title={Grammar Lens: -\ml{ഉം} + \ml{ആയിരുന്നു}}]
Put \textbf{\ml{ആയിരുന്നു}} after the future form, and the two make \textquotedblleft{}used to\textquotedblright{}:

\begin{quote}
\ml{പോകും} + \ml{ആയിരുന്നു} $\to$ \textbf{\ml{പോകുമായിരുന്നു}} — \emph{pōkumāyirunnu} — used to go
\end{quote}

\begin{quote}
\textbf{\ml{കുട്ടിക്കാലത്ത്} \ml{ഞാൻ} \ml{എല്ലാ} \ml{ദിവസവും} \ml{കുളത്തിൽ} \ml{കുളിക്കുമായിരുന്നു}.}
\end{quote}

\begin{quote}
\emph{kuṭṭikkālattŭ ñān ellā divasavuṁ kuḷattil kuḷikkumāyirunnu.} — \textquotedblleft{}In my childhood I used to bathe in the pond every day.\textquotedblright{}
\end{quote}

\begin{quote}
\textbf{\ml{പണ്ട്} \ml{മുത്തശ്ശി} \ml{രാവിലെ} \ml{കടപ്പുറത്ത്} \ml{നടക്കുമായിരുന്നു}.}
\end{quote}

\begin{quote}
\emph{paṇṭŭ muttaśśi rāvile kaṭappuṟattŭ naṭakkumāyirunnu.} — \textquotedblleft{}Long ago Grandmother used to walk on the beach in the morning.\textquotedblright{}
\end{quote}

\textbf{\ml{എല്ലാ} \ml{ദിവസവും}} is \textquotedblleft{}every day\textquotedblright{}. \textbf{\ml{കുട്ടിക്കാലത്ത്}}, \textquotedblleft{}in childhood\textquotedblright{}, and \textbf{\ml{പണ്ട്}}, \textquotedblleft{}long ago\textquotedblright{}, are the words that usually stand in front of a habit.
\end{grammarlens}

\begin{grammarlens}[title={Grammar Lens: three pasts side by side}]
\noindent
\begin{tabularx}{\linewidth}{@{}>{\raggedright\arraybackslash}X>{\raggedright\arraybackslash}X@{}}
\toprule
\textbf{Malayalam} & \textbf{what it says} \\
\midrule
\textbf{\ml{ഞാൻ} \ml{നടന്നു}} (\emph{ñān naṭannu}) & I walked: once, and it is done \\
\textbf{\ml{ഞാൻ} \ml{നടക്കുകയായിരുന്നു}} (\emph{ñān naṭakkukayāyirunnu}) & I was walking: the scene, still running \\
\textbf{\ml{ഞാൻ} \ml{നടക്കുമായിരുന്നു}} (\emph{ñān naṭakkumāyirunnu}) & I used to walk: again and again \\
\bottomrule
\end{tabularx}

Hold on to the last one. In chapter 348 the same \textbf{-\ml{ഉമായിരുന്നു}} comes back with a second job: \textquotedblleft{}would have\textquotedblright{}.
\end{grammarlens}

\subsection*{Your turn}
Say these aloud:

\begin{itemize}
\raggedright
  \item \textquotedblleft{}I used to go\textquotedblright{}, \textquotedblleft{}I used to read\textquotedblright{}, \textquotedblleft{}I used to walk\textquotedblright{}
  \item \emph{kuṭṭikkālattŭ ñān ellā divasavuṁ kuḷattil kuḷikkumāyirunnu}
  \item one thing your grandmother used to do, with \emph{paṇṭŭ} in front
  \item I walked; I was walking; I used to walk
\end{itemize}

\begin{tcolorbox}[breakable,colback=teal!4,colframe=teal!35!black,title={Before you move on}]
\textquotedblleft{}Used to go\textquotedblright{}? (\emph{\textbf{pōkumāyirunnu}}: the future \textbf{\ml{പോകും}} with \textbf{\ml{ആയിരുന്നു}}.) \textquotedblleft{}I was walking\textquotedblright{}? (\emph{\textbf{ñān naṭakkukayāyirunnu.}}) \textquotedblleft{}I walked\textquotedblright{}, once? (\emph{\textbf{ñān naṭannu.}})
\end{tcolorbox}
\section[(review)]{Review — the past three ways}

\label{lesson:ML-R339-the-past-three-ways}



\begin{quote}
Say \textquotedblleft{}he was at home yesterday\textquotedblright{}. (\emph{avan innale vīṭṭil āyirunnu.}) Then \textquotedblleft{}there was a small shop in the village\textquotedblright{}. (\emph{grāmattil oru ceṟiya kaṭa uṇṭāyirunnu.})
\end{quote}

\subsection*{Your turn}
Say these aloud:

\begin{itemize}
\raggedright
  \item \textquotedblleft{}That was a good film\textquotedblright{} — \emph{atŭ oru nalla sinima āyirunnu}
  \item \textquotedblleft{}Grandmother had a big house\textquotedblright{}, then \textquotedblleft{}there was no phone there\textquotedblright{}
  \item \textquotedblleft{}I was reading a book. Then Mother called.\textquotedblright{}
  \item \textquotedblleft{}We were walking on the beach. Suddenly it rained.\textquotedblright{}
  \item \textquotedblleft{}In my childhood I used to bathe in the pond every day.\textquotedblright{}
  \item a scene from yesterday evening, in the -ukayāyirunnu past, and one thing that cut into it
\end{itemize}

\begin{tcolorbox}[breakable,colback=teal!4,colframe=teal!35!black,title={Before you move on}]
\textquotedblleft{}Was\textquotedblright{}? (\textbf{\ml{ആയിരുന്നു}}.) \textquotedblleft{}There was\textquotedblright{}, \textquotedblleft{}had\textquotedblright{}? (\textbf{\ml{ഉണ്ടായിരുന്നു}}.) \textquotedblleft{}Was reading\textquotedblright{}? (\textbf{\ml{വായിക്കുകയായിരുന്നു}}.) Which past paints the scene, and which drops in the event? (The \emph{-ukayāyirunnu} past is the scene; the plain past, \emph{\textbf{viḷiccu}}, \emph{\textbf{peytu}}, is the event.) \textquotedblleft{}Used to walk\textquotedblright{}? (\textbf{\ml{നടക്കുമായിരുന്നു}}, \emph{naṭakkumāyirunnu}.)
\end{tcolorbox}
